\chapter{Objetivos y alcance}

\section{Objetivo general}
\pendiente{Reformular sobre el aprobado; corregir la norma: MPEG-2 Video es
ISO/IEC 13818-2 (el anteproyecto citaba ISO/IEC 23002-1, que es la de
precisión de la IDCT).}

\section{Objetivos específicos}
\pendiente{Lista del anteproyecto, marcando cuáles se redefinieron.}

\section{Redefinición del alcance}
\label{sec:redefinicion-alcance}
% Reglamento 2.3: se admiten las acciones "que la experiencia adquirida
% indique como oportunas y necesarias". Criterio de evaluación 3.6 (e).
\subsection{Sistema operativo: Linux en lugar de un RTOS propio}
\pendiente{Acuerdo verbal con el tutor, Ing.~Ricardo Maldonado, a comienzos
de 2026 (febrero): acotar el alcance para lograr un plazo de finalización
realista. Qué objetivos del anteproyecto reemplaza (``Diseñar el sistema
operativo'', ``Implementar el sistema operativo'', etapas 6 y 7) y a qué se
destinó ese esfuerzo (driver de kernel, protocolo, producto). La
justificación técnica va en el capítulo~\ref{cap:firmware}.}

\subsection{Salida de video}
\pendiente{La PolarFire SoC Discovery Kit no tiene conector de video: la
cadena RTL de visualización (yuv2rgb, syncgen) queda instanciada pero sin
verificar, y la conversión YUV$\to$RGB se hace en el cliente.}

\section{Alcance final}
\pendiente{Qué entra y qué no (p.\,ej.\ matrices de cuantización cargadas
como limitación conocida, si sigue abierta al entregar).}
